\documentclass{beamer}
\usepackage{amsmath}
\setbeameroption{show notes on second screen=right}
\setbeamertemplate{note page}[plain]

\beamertemplatenavigationsymbolsempty

\begin{document}

\begin{frame}
\frametitle{}

$\frac{15}{4} \div \frac{2}{5}$
\note[item]{\alert<1>{Find 15 over 4 divided by 2 over 5.}}

\vskip4pt

\pause $=\frac{15}{4} \times \frac{5}{2}$
\note[item]{\alert<2>{To divide fractions, we multiply by the reciprocal. We only reciprocate the second fraction, and turn division into multiplication at the moment recipcrocation happens.}}
\note[item]{\alert<3>{So dividing by 2 over 5 turns into multiplying by 5 over 2. Our problem is now 15 over 4 times 5 over 2.}}

\vskip4pt

\pause
\pause $=\frac{15 \times 5}{4 \times 2}$ 
\note[item]{\alert<4>{To multiply fractions, we multiply straight across (whether we have a common denominator or not). So the numerator is 15 times 5, and the denominator is 4 times 2.}}

\vskip4pt

\pause $=\frac{75}{8}$
\note[item]{\alert<5>{Simplifying the numerator and denominator separately, we have 75 over 8.}}
\note[item]{\alert<6>{Since 75 and 8 do not share any common factors, this is as simplified as it gets as an improper fraction.}}
\note[item]{\alert<7>{In most situations in algebra, we leave our answer like this.}}
\note[item]{\alert<8>{If we need to write this as a mixed fraction, 75 divided by 8 gives us 9 with a remainder of 3, so we would have 9 and 3 over 8. Both are correct: 75 over 8 is correct, and 9 and 3 over 8 is also correct.}}

\pause 
\pause 
\pause 
$=9\frac{3}{8}$

\end{frame}





%% Blank frames at the end to prevent weird shifts.
\begin{frame}
\end{frame}
\begin{frame}
\end{frame}
\begin{frame}
\end{frame}
\begin{frame}
\end{frame}
\begin{frame}
\end{frame}
\begin{frame}
\end{frame}
\begin{frame}
\end{frame}

\end{document}
